\section{The result and its foundational meaning}

The ordinary implicit function theorem starts with a \(C^1\) map
\(F:\R^n\times\R^m\to\R^m\), a zero \((a,b)\), and an invertible
\(y\)-Jacobian at that zero. It obtains a local graph \(y=g(x)\) and the
formula
\[
 Dg(x)=-D_yF(x,g(x))^{-1}D_xF(x,g(x)).
\]
Our result preserves these conclusions for quantitatively presented
finite-dimensional data. It additionally specifies how to compute the
neighborhood, the graph, the derivative, and all relevant error bounds.
The contraction argument is classical in origin; the contribution here
is a complete finite-observation organization adapted to
quantifier-free induction, rather than a claim of a new classical
implicit function theorem. Constructive implicit function theorems and
their quantitative uniqueness principles have an existing literature,
including Diener--Schuster \cite{diener2011,diener2009}.

\subsection{The base theory}

We use first-order intuitionistic primitive recursive arithmetic:
intuitionistic predicate logic with equality, the natural numbers
\(0,S\), a function symbol for every fixed primitive recursive
function, their defining equations, the usual elementary arithmetic
axioms, and induction
\[
 P(0,\vec u)\ \wedge\
 \forall j\bigl(P(j,\vec u)\to P(j+1,\vec u)\bigr)
 \ \longrightarrow\ \forall j\,P(j,\vec u)
 \tag{QF-Ind}
\]
only when \(P\) is quantifier-free. Equivalent finite-signature
presentations can be used through a previously verified interpretation;
we do not need such an interpretation for the mathematical proof.
An order comparison between rational codes is decidable and can be
represented by a primitive recursive Boolean function.

All dimensions in the theorem are fixed positive integers. A finite
tuple of rational numbers, a matrix, a finite trace, a finite sum, and
a bounded maximum are coded by natural numbers. The coding and the
translation of the argument into (QF-Ind) are made explicit in
Section~\ref{sec:arith}.

\subsection{What a real function in iPRA means}

A primitive recursive real vector is presented by a fixed primitive
recursive map \(u:\N\to\Q^d\) satisfying
\begin{equation}\label{eq:regular}
 \norm{u(k)-u(l)}\le 2^{-k}+2^{-l}\qquad(k,l\in\N).
\end{equation}
This is a first-order universal statement about numerical codes.
Writing a real vector \(u\), its norm, or an inequality about it is
convenient notation for corresponding statements at every rational
observation precision. No set of all real numbers is introduced into
the object theory.

A function on a rational box is represented by a fixed primitive
recursive evaluator at rational points, together with uniform
continuity data. Its action on any fixed primitive recursive point
name is then another primitive recursive program, obtained by
composition and an explicit precision shift. Uniform statements about
the evaluator use ordinary numerical quantifiers. Statements about
arbitrary named input points are schemas over their fixed names.

This distinction is essential. There is no primitive recursive universal
evaluator for all primitive recursive program codes. If such a total
evaluator \(U(e,x)\) existed, the primitive recursive function
\(x\mapsto U(x,x)+1\) would contradict universality at its own code.
We fix the input evaluators before taking primitive recursion. A
uniform numerical family of inputs is also allowed when its evaluator
is already primitive recursive jointly in the numerical family
parameter.

\subsection{The certificates that replace hidden existence assertions}

The input contains a neighborhood radius, positive conditioning data,
and primitive recursive moduli for a Taylor remainder and continuity
of the derivative. These are substantive hypotheses. Merely writing
that \(F\) is \(C^1\) without specifying the presentation, moduli, and
their arithmetic proofs is not enough for an iPRA construction.
Similarly, a negation \(\neg(\det J=0)\) is not silently turned into a
positive numerical lower bound for \(|\det J|\).

For polynomial maps with rational coefficients, these certificates
come from finite polynomial identities and explicit box estimates.
The same interface can be discharged for other functions with
primitive recursive evaluation and certified local calculus. The
theorem does not assert that an arbitrary computable function has a
primitive recursive evaluator or a primitive recursive modulus.

\subsection{Norms and a bounded logarithmic shift}

Vectors carry the maximum norm; matrices carry its induced row-sum
norm:
\[
 \norm{z}=\max_i|z_i|,\qquad
 \norm{B}=\max_i\sum_j|B_{ij}|.
\]
The product \((x,y)\) also carries the maximum norm. Thus
\(\norm{Bz}\le\norm B\norm z\) and
\(\norm{BC}\le\norm B\norm C\). For rational matrices these statements
are finite triangle inequalities.

For a nonnegative rational \(t\), let \(\lgup(t)\) be the least
integer \(j\) in the bounded range \(0\le j\le\lceil t\rceil+1\)
such that \(t\le2^j\). This is a primitive recursive bounded search;
the range contains a witness by elementary induction. No real
logarithm is used. In particular,
\[
 t\,2^{-\lgup(t)}\le1.
\]
Every logarithmic-looking precision constant below means this bounded
operation.
